% =====================================================================
\lezione{14/10/2026}{RNN in PyTorch, CNN e graph neural network}
  {Prof.~Danilo Giordano}
  {lunghezze diverse e padding dinamico; \code{collate\_fn} e packed sequence; \code{out} e
   \code{hidden}; notebook pollution e FordA; CNN (convoluzione, pooling, reti famose, STAMINA); GNN
   (permutation invariance ed equivariance, message passing, GCN, GraphSAGE, GAT, over-smoothing,
   task, split transductive e inductive, limiti)}
\label{lez:2026-10-14}
% =====================================================================

\begin{nota}[Fonte]
Trascrizione automatica della registrazione, revisionata. Codice, formule e numeri sono ripresi dalle
slide \emph{3-deep-learning-methodologies} (pp.~72--123) e dai notebook
\emph{RNN-Pollution-Multivariate} e \emph{RNN-LSTM-FordA}.
\end{nota}

\section{Sintesi}
Lezione lunga che chiude la parte sulle architetture per la classificazione. Prima si completa la
pratica delle RNN: come gestire serie di lunghezza diversa con il padding dinamico nel mini-batch e le
\en{packed sequence}, e due notebook (regressione sull'inquinamento, classificazione FordA). Poi una
breve introduzione alle \textbf{CNN}, utile anche per introdurre il concetto di convoluzione, e infine
la teoria delle \textbf{graph neural network}. Approfondimento nei capitoli~\ref{cap:rnn} e
\ref{cap:cnn-gnn}.

\section{Avvisi pratici}
\begin{itemize}
  \item Il \textbf{Lab~2} si terrà nei due mercoledì successivi (21 e 28 ottobre). Per liberare il
    mercoledì, la settimana prossima ci sarà lezione \textbf{lunedì e martedì} (19 e 20 ottobre);
    l'avviso sarà pubblicato sul portale.
  \item Nel Lab~2 le sequenze di API call descrivono un goodware o un malware: servirà la
    rappresentazione dell'\textbf{intera serie}; la parte GNN dei lab fa solo graph classification.
  \item Le GNN in PyTorch si vedranno lunedì (\rif{2026-10-19}{19/10}).
\end{itemize}

\section{Contenuti}

\subsection{RNN in pratica: lunghezze diverse}
Una rete è moltiplicazione di matrici: una serie più corta lascia un ``buco''. Tagliare tutto alla
lunghezza minima butta informazione; paddare con valori inventati (es.\ la media) fa credere alla rete
che la serie finisca così. Soluzione: un \code{Dataset} personalizzato e una \code{collate\_fn},
passata al \code{DataLoader}, che padda le serie \textbf{solo fino alla lunghezza massima del
batch} e restituisce anche le lunghezze originali. Basta paddare nel batch perché tutti i passi
condividono gli stessi pesi. Nel modello, \code{pack\_padded\_sequence} dice alla RNN quali istanti
sono padding, così vengono saltati.

\begin{esame}
Il padding serve solo alla matematica (matrice completa): \textbf{non deve entrare} nello stato finale
né nella backpropagation, altrimenti la rete si addestra su dati inventati. Il padding vale 0, che
potrebbe essere un valore reale: per questo la RNN deve sapere dove inizia.
\end{esame}

\subsection{\code{out} e \code{hidden}}
\code{hidden} contiene lo stato finale di ogni layer (\code{hidden[-1]} è quello dell'ultimo layer, la
rappresentazione più astratta) e si usa per i task a livello di \textbf{sequenza}, come la
classificazione; \code{out} contiene gli stati dell'ultimo layer a ogni passo e serve per i task a
livello di \textbf{passo} (previsione del valore successivo, many-to-many). Esempio della scalata:
\code{hidden} riassume tutto il viaggio, \code{out} lo descrive passo per passo. Varianti:
bidirezionale (head di dimensione \code{2*hidden\_size}, concatenando le due direzioni) e LSTM
(restituisce hidden e cell state; per classificare si usa l'hidden).

\subsection{I due notebook}
\begin{itemize}
  \item \textbf{Pollution} (regressione): guardare i dati prima di usarli (direzione del vento come
    stringa, valori sospetti come vento a 600~km/h); niente split casuale che rompe le serie; scaling
    fittato sul training; RNN + layer lineare senza attivazione finale, loss MSE. Il training loop è
    identico a quello della FFNN; le loss delle RNN sono più rumorose; manca il plateau; l'MSE va letto
    rispetto alla scala dei valori.
  \item \textbf{FordA} (classificazione di serie univariate): per mostrare il padding si accorciano a
    480 le prime 100 serie (le altre sono lunghe 500). Risultato: accuracy intorno al 51\% su un dataset
    bilanciato, predicendo sempre la stessa classe, con RNN, LSTM e Bi-RNN. Non perché le RNN non
    funzionino: head troppo semplice, hidden size piccola, training breve.
\end{itemize}

\subsection{Convolutional neural network}
Usi in cybersecurity: malware come immagine, siti clone da screenshot, intrusion detection su serie
(anche CNN-LSTM), deepfake. Struttura: layer convoluzionali che estraggono feature in modo gerarchico
(bordi, forme, oggetti) e una FFNN finale. Il layer è \textbf{convoluzione} (filtri scorrevoli con
pesi addestrabili; $N$ filtri producono $N$ strati di profondità; padding ai bordi), \textbf{ReLU} e
\textbf{pooling} (es.\ max pooling $2\times2$: una matrice $4\times4$ diventa $2\times2$). Reti
famose: LeNet-5 (1998), AlexNet (2012, ReLU e dropout), ResNet-50 (2015, skip connection).
\textbf{STAMINA} (Intel): byte del binario come pixel in scala di grigi, immagine $224\times224$,
ResNet-50 pre-addestrata con fine-tuning, accuracy oltre il 95\%.

\subsection{Graph neural network}
Molti dati sono grafi (botnet, frodi, grafi di API call). Una FFNN su matrice di adiacenza e feature
dipende dall'ordine dei nodi, non scala e non gestisce grafi di dimensione diversa; CNN e RNN
presuppongono una struttura fissa. Si generalizza la convoluzione: \textbf{un nodo è descritto dai
suoi vicini}. Due proprietà: \textbf{permutation invariance} a livello di grafo,
$f(A_1,X_1)=f(A_2,X_2)$, e \textbf{permutation equivariance} a livello di nodo,
$Pf(A,X)=f(PAP^T,PX)$. Con il \textbf{message passing} ogni nodo aggrega i messaggi dei vicini
(che a loro volta hanno aggregato i propri) con pesi \textbf{condivisi tra tutti i nodi}; un layer è
messaggio + aggregazione indipendente dall'ordine (somma, media, max) + il messaggio del nodo stesso +
non linearità.

\begin{esame}
Saper spiegare \textbf{perché} una GNN è invariante/equivariante (pesi condivisi e aggregazione che non
dipende dall'ordine) e perché è \textbf{induttiva} (gli stessi parametri descrivono nodi mai visti).
Il docente insiste anche sull'\textbf{attenzione} di GAT, che tornerà nella parte di NLP.
\end{esame}

Tre architetture: \textbf{GCN} (tutti i vicini, normalizzati per il grado), \textbf{GraphSAGE}
(campiona un numero fisso di vicini, aggregazione in due stadi: adatta a grafi grandi), \textbf{GAT}
(pesa ogni vicino con un coefficiente di attenzione appreso). Con troppi layer i \en{receptive field}
si sovrappongono e gli embedding diventano tutti simili (\textbf{over-smoothing}): scegliere il
numero di layer in base al receptive field necessario, aumentare l'espressività dentro il layer,
aggiungere MLP prima e dopo, skip connection. Task: node classification, link prediction, graph
classification. Split \textbf{transductive} (un solo grafo, si dividono le label) o
\textbf{inductive} (grafi diversi negli split, il setting dei lab). Limiti: over-smoothing,
scalabilità, grafi eterogenei e dinamici.

\section{Domande aperte e punti da verificare}
\begin{verifica}
\begin{itemize}
  \item Date del Lab~2 e lezione del martedì: ricavate da ``\en{next week}'' e ``\en{the next one and
    the one after}'' detti in aula; confermare con l'avviso sul portale.
  \item Nel notebook pollution si normalizza anche la label; il docente osserva che probabilmente non
    andrebbe fatto.
\end{itemize}
\end{verifica}
